% !TEX root = ../../main.tex

\section{Artefactos de Validación}

La validación deja artefactos reproducibles para inspeccionar el comportamiento del compilador y respaldar los resultados. Los experimentos viven bajo \texttt{experiments/}, mientras que los resultados se guardan bajo \texttt{tests/} siguiendo una estructura paralela por mónada, familia, caso y variante.

Los archivos más relevantes son:

\begin{center}
\begin{tabular}{ll}
\hline
Artefacto & Función \\
\hline
\texttt{summary.log} & Resume costos y cantidad de permutaciones. \\
\texttt{semantic-validation.log} & Registra compilaciones, ejecuciones y comparación final. \\
\texttt{logs/raw.log} & Guarda la traza cruda del Renamer. \\
\texttt{logs/graph.log} & Presenta el grafo RAW en forma legible. \\
\texttt{logs/optimal-reorder.log} & Muestra el candidato elegido automáticamente. \\
\texttt{logs/original\_ado.log} & Muestra \texttt{ApplicativeDo} normal, sin reordenamiento. \\
\texttt{logs/permutation\_i.log} & Muestra el candidato \texttt{i} forzado por flag. \\
\hline
\end{tabular}
\end{center}

Estos artefactos cumplen dos roles. Primero, permiten depurar la extensión observando el grafo, las permutaciones y los árboles seleccionados. Segundo, entregan la base empírica para comparar costos y salidas entre variantes.
